
% ============================================================================
\clearpage
\section*{v4.10.0-r16 Source-Only: Status Notes on the Helicity Axis, the Invariant Doublet and the Frame Signs}
\addcontentsline{toc}{section}{v4.10.0-r16 Source-Only: Status Notes on the Helicity Axis, the Invariant Doublet and the Frame Signs}
\label{sec:v4100-r16-helicity-invariant-doublet}
% ============================================================================

\paragraph{Purpose.}
This layer records the OP3 explorations of 15 November 2026 that followed v4.10.0-r15, and an earlier exploration of
7 November 2026 that had remained on a local branch.  It restricts the reading of v4.10.0-r15.5 on the states of
\(\mathbb C\otimes\OO\), places the Koide phase on the helicity axis of the field doublet \((E,B)\) of v4.9.0-r04 and r05, moves
the octant onto a Lorentz-invariant doublet, identifies the three axes with three signs of a field that the frame maps
preserve one at a time, excludes a candidate for the missing threefold cycle, and shows that a condition of
non-orthogonality makes the cycle unnecessary.  All items are explorations without
pre-registration and without comparison with data; the Koide value \(K=\tfrac23\), the phases and most outcomes were fixed
analytically before the checks.  The notes are \texttt{notes/\allowbreak research/\allowbreak v4\_10\_0\_r16/} and the
folders cited there.  Earlier layers are not rewritten.

\noindent\textbf{Source-only status.}
This v4.10.0-r16 layer is source-only.  It does not build a PDF, does not assign a new DOI,
and does not authorize a public release.  The v5.1.0 Version DOI
\href{https://doi.org/10.5281/zenodo.24237112}{10.5281/zenodo.24237112} and the Concept DOI
\href{https://doi.org/10.5281/zenodo.19210728}{10.5281/zenodo.19210728} are unchanged by it.

\subsection*{v4.10.0-r16 status notes}
\label{subsec:v4100-r16-status}

\paragraph{Theoremlet v4.10.0-r16.1 (restriction of r15.5 on \(\mathbb C\otimes\OO\)).}
With the operators of v4.9.0-r51 on \(\mathbb C\otimes\OO\): the tachyon map \(R\) is multiplication by \(i\), a scalar, so every
state is an eigenstate; \(T\) and \(J_B\) are antilinear and have the same eigenrays, the real states up to a phase, with mutual
overlap \(1\); these states carry equal weight in \(P_{(\ell,r)}\) and \(P_{(-\ell,-r)}\) and have \(\langle Q\rangle=\tfrac12\).  The three
mutually unbiased bases of r15.5 exist only on the complexified amplitude plane \(z\), not among the matter states.  Read on
\(\mathbb C\otimes\OO\), two of the three ``generations'' of r15.5 would have no definite charge.  This restricts the reading of
r15.5; its computation (a)--(c) stands on the plane \(z\).

\paragraph{Theoremlet v4.10.0-r16.2 (the Koide phase needs the helicity axis).}
Let \(b=(b_1,b_2,b_y)\) be the Stokes vector of the doublet \((E,B)\): \(b_1\), \(b_2\) as in v4.9.0-r04.2 and \(b_y\) the helicity
component of v4.9.0-r05.2(b), \(-2\operatorname{Im}(\hat E^*\cdot\hat B)\) relative to the trace.
\begin{enumerate}[leftmargin=2em]
\item[(a)] The tachyon map \((E,B)\mapsto(-B,E)\) is the duality rotation by \(90^\circ\), a half turn about the helicity axis: it
reverses \(b_1\), \(b_2\) and keeps \(b_y\).  \(T\colon(E,B)\mapsto(E,-B)\) reverses \(b_2\) and \(b_y\); \(J_B\colon(E,B)\mapsto(B,E)\)
reverses \(b_1\) and \(b_y\).  This is the pattern of v4.9.0-r51.1(a) for chirality.
\item[(b)] Under the map \(\varphi\) of v4.9.0-r51.1(c) the \(\sigma_y\) eigenvectors \((1,\pm i)\) go to \(1\pm ie_7\), the
eigenvectors of \(iL_{e_7}\) for \(\pm1\): compensator and partner port.
\item[(c)] Three states with \(b_y=0\) have Bargmann phase \(0\) or \(\pi\).  For every octant triple (three orthonormal Bloch
vectors) \(\sum_k b_{y,k}^2=1\).  Real fields alone cannot carry the Koide phase; it needs complex amplitudes.
\item[(d)] For general complex amplitudes \(b_y\) changes under boosts; only the splitting into the helicity parts
\(F_\pm=E\pm iB\) is preserved.
\end{enumerate}

\paragraph{Theoremlet v4.10.0-r16.3 (a Lorentz-invariant doublet).}
Put \(s_\pm=F_\pm\cdot F_\pm\) (bilinear).  A boost acts on \(F_+\) by \(M\in SO(3,\mathbb C)\) and on \(F_-\) by \(\bar M\), so \(s_\pm\) are
invariant (v4.8.0-r06.3), and \(\chi=(\sqrt{s_+},\sqrt{s_-})\) is a doublet whose Bloch sphere does not depend on the observer.
\begin{enumerate}[leftmargin=2em]
\item[(a)] Real fields lie on the equator at angle \(-\arg s\), the direction in the disk of r04; self-dual and anti-self-dual
non-null fields lie at the poles; for plane waves \(s_+=s_-=0\) and light has no point.
\item[(b)] \(R\) turns the equator by \(\pi\), \(T\) reflects \(\arg s\mapsto-\arg s\), \(J_B\) reflects \(\arg s\mapsto\pi-\arg s\), as in
r04.2(b); \(R\) keeps the helicity sign, \(T\) and \(J_B\) reverse it.
\item[(c)] The roots carry signs.  For real fields \(\sqrt{s_-}=\overline{\sqrt{s_+}}\) is a natural choice; the other choice gives
the image under \(R\).  For general complex fields a point is fixed only up to \(R\); at the poles the choice is void.
\item[(d)] A matter field (\(s=1\)), a field of the photon side (\(s=i\)) and a self-dual field (\(E=iB\)) give pairwise overlap
\(\tfrac12\), \(K=\tfrac23\) and Bargmann phase \(\pm\pi/4\), and all three rays are unchanged under boosts.
\end{enumerate}

\paragraph{Theoremlet v4.10.0-r16.4 (three signs of a field and the frame maps).}
The axes of r16.3 carry \(\operatorname{sign}\operatorname{Re}s\) (matter against tachyon), \(\operatorname{sign}\operatorname{Im}s\) (the side
of v4.9.0-r02.1) and the helicity sign \(\operatorname{sign}(|s_+|-|s_-|)\).
\begin{enumerate}[leftmargin=2em]
\item[(a)] Each of \(R\), \(T\), \(J_B\) reverses two signs and keeps exactly one: \(R\) the helicity, \(T\) \(\operatorname{Re}s\), \(J_B\)
\(\operatorname{Im}s\).  On the signs the three maps form the Klein four-group.
\item[(b)] The poles of each axis are eigenstates of the map that keeps its sign (matter and tachyon field for \(T\), the two
side fields for \(J_B\), the self-dual and anti-self-dual fields for \(R\)), and of no other.
\item[(c)] A linear map of the index plane, \((F_+,F_-)\mapsto(aF_++bF_-,cF_++dF_-)\), gives
\(s_+'=a^2s_++2ab\,F_+\cdot F_-+b^2s_-\).  The mixed term is not Lorentz invariant, so only diagonal and exchanging maps act on
\(\chi\); they keep or exchange the poles.  No linear field map cycles the three axes.
\end{enumerate}
Reading: if each of three readout states is required to be an eigenstate of a different frame map, the overlap \(\tfrac12\),
\(K=\tfrac23\) and the phase \(\pm\pi/4\) follow and are boost invariant; ``a complete set'' of r15.4 is no longer needed, since
there are exactly three frame maps.  The rule ``one readout state per frame map'' remains assumed.  In the Gram reading the
square roots of the masses are the eigenvalues of the Gram matrix; the three states are readout states, not mass eigenstates.

\paragraph{Theoremlet v4.10.0-r16.5 (the Hurwitz cycle does not supply the threefold cycle).}
On the neutral qubit \(N\) of v4.10.0-r13 (under CT-003) the tachyon map is the scalar \(i\), and complex conjugation maps
\(N\) onto the charge-conjugate space \(\bar N\), orthogonal to \(N\).  The frame maps act on \(N\) differently from \(\chi\), so no
identification of \(N\) with \(\chi\) compatible with \(R\), \(T\), \(J_B\) exists.  The threefold cycle of the Hurwitz units of r13.2
permutes the Pauli axes of the family \(\mathfrak{su}(2)\) of r13.4, not the frame maps.  On the operator level the Clifford
element of order three cycles \(T\sim\sigma_z\), \(J_B\sim\sigma_x\), \(R\sim\sigma_y\) by conjugation, but as a field map it is not
diagonal and does not act on \(\chi\).

\paragraph{Theoremlet v4.10.0-r16.6 (``one object'': orbits of \(SO(3,\mathbb C)\times U(1)\); exploration of 7 November 2026).}
\(G=SO(3,\mathbb C)\times U(1)_{\rm duality}\) acts on \(z=E+iB\) transitively on every shell \(|s|=r>0\) (constructive check on 2000
pairs, residual at most \(1.2\cdot10^{-14}\); different shells cannot be joined).  Its orbits are the shells, the punctured null
cone and \(\{0\}\), and \(s=0\) is the only fixed point of duality.  Matter and tachyon lie on one orbit, ``one object'' in the sense
of an active symmetry of the free carrier; light lies on another and is joined to them only through the common carrier or as
a limit.  The identifications of the core fall into kinds: symmetries and changes of observer keep information; the
compensator, the mixture \(P=I/2\) and the Hopf reading of r61 lose or mix it.  The counting rule per frame belongs to none of
these kinds and cannot come from \(G\).  The exploration was kept on a local branch until 15 November 2026.

\paragraph{Theoremlet v4.10.0-r16.7 (readouts of one object: at most three, and they form an octant).}
On \(\chi\) the frame maps act as \(R=\sigma_z\), \(T=\sigma_x\), \(J_B=\sigma_y\) up to phase (checked against the field action at the six
poles and on 200 real fields, deviation \(4\cdot10^{-16}\)); each has exactly two eigenrays, and together these are the six poles of
r16.4(b).  Among the subsets of the six poles, 26 contain no orthogonal pair; the largest have three elements, and there are
exactly eight of them.  All eight have pairwise overlap \(\tfrac12\), \(K=\tfrac23\) and Bargmann phase \(\pm\pi/4\), and the frame maps
split them into two orbits of four, one for each sign of the phase.
Reading: require (1) that a readout is an eigenstate of a frame map, (2) that two readouts of one object are not orthogonal,
since orthogonal states exclude each other and would be different objects (meaning (4) of ``one object'' in r16.6), and (3) that
all such readouts occur.  Then the number three, the overlap \(\tfrac12\) and the phase \(\pi/4\) follow together, and the triple is
unique up to frame maps and orientation.  No threefold cycle is needed: two readouts of the same frame map would be
orthogonal.  The rule of r16.4 is thereby reduced to (2) and (3); (1), (3), the sign of the phase, the bend to \(\varphi=\tfrac23\) and a
charged carrier remain open.  Not blind: the targets three and \(\tfrac23\) were known, and the conditions were stated afterwards.

\subsection*{v4.10.0-r16 status in the three frames}
\label{subsec:v4100-r16-frames}

\begin{itemize}[leftmargin=2em]
\item M, T: r16.3 and r16.4 are invariant under boosts; \(R\) exchanges the matter and tachyon signs and keeps the helicity,
so the Koide data on \(\chi\) are the same in M and T.
\item P: light has no point on \(\chi\) (r16.3(a)); by r16.6 it is the fixed point of duality, not on the orbit of M and T.
\end{itemize}

\subsection*{v4.10.0-r16 register entries}
\label{subsec:v4100-r16-register}

\begin{center}
\small
\begin{tabular}{@{}>{\raggedright\arraybackslash}p{0.34\textwidth}>{\raggedright\arraybackslash}p{0.58\textwidth}@{}}
\toprule
statement & status after v4.10.0-r16\\
\midrule
one generation per frame map on \(\mathbb C\otimes\OO\) (r15.5) & does not hold there; only on the plane \(z\) (r16.1)\\
Koide phase and helicity & phase needs the helicity axis; frame maps act on helicity as on chirality (r16.2)\\
invariant doublet \(\chi\) & boost-invariant octant of three field types; root branch fixed up to \(R\) (r16.3)\\
three field signs & each frame map keeps one; overlap \(\tfrac12\) conditional on one readout state per frame map (r16.4)\\
threefold cycle & not from linear field maps, not from the Hurwitz units on \(N\) (r16.4(c), r16.5); not needed under r16.7\\
readouts of one object & at most three non-orthogonal frame eigenstates; they form an octant, unique up to frame maps and orientation (r16.7)\\
one object (k220) & M and T one orbit, light the fixed point; counting rule not from \(G\) (r16.6)\\
CT-001 to CT-004 & unchanged\\
H100 / H102 & pending / pending\\
\bottomrule
\end{tabular}
\end{center}

\subsection*{v4.10.0-r16 ledger and guard update}
\label{subsec:v4100-r16-ledger-guard}

\[
R_{\rm chk}^{v4100\text{-}r16}=(1,1,0)
\]
for (status notes recorded; formulas checked in the cited scripts; the eigenstate condition and maximality of r16.7, the core
size and a charged carrier not derived).

The layer spends no guard:
\[
R_{\rm glob}^{(v4100\text{-}r16,+)}=(0,1,1,1,1),
\qquad
\Gamma_{\rm guard}^{v4100\text{-}r16}=(1,1,1,1,1),
\qquad
\beta_4=1/256\ \text{unspent}.
\]
OP~1, OP~3 and OP~4 remain open.  The layer does not imply exact finite \(C5\) zero closure, exact all-readout terminality,
\(C6/T2\), a global photon-clock/RH theorem, or global all-readout closure.
